\section{Scope and Feature Overview}
\label{sec:scope}

\subsection{Included Scope}
OliveTrack is privately delivered to Pepe's parents for one client business, with
one initial owner account prepared by a Technical Administrator. It supports:
account access, delegated management, properties and optional parcels, campaigns,
workers, optional planning, attendance and actual work, worker documents and signing,
deliveries, cooperative results and campaign indicators.

Table~\ref{tab:feature-overview} maps these areas to their use cases. Detailed
behaviour, permissions and correction conditions are specified in
Sections~\ref{sec:use-cases}--\ref{sec:business-rules}. Supported field operations
may continue offline; connectivity is defined per operation. Uploaded identity and
work-permit evidence is not issued or verified by document generation.

\subsection{Excluded Scope}

Product scope describes supported system capabilities; project scope also
includes the work needed to specify, build, verify and deliver them. The scope
below is the current working baseline for a single final delivery under
\decisionref{Q-SCOPE-01}. Individual pending scope choices are resolved explicitly
before changing that baseline.

The current proposal excludes:

\begin{itemize}
    \item Payroll execution, wage calculation,
    worker remuneration rates and payment processing.
    \item Invoicing, full accounting
    and commercial olive oil sales.
    \item Application-calculated price estimates
    for deliveries.
    \item Worker self-service accounts
    for attendance or business management.
    \item Direct cooperative or laboratory
    integrations.
    \item Automatic extraction of values
    from ticket images or documents.
    \item Independent transport-trip management
    beyond the delivery-associated transport record.
    \item Management of agricultural activities
    unrelated to the agreed olive-harvest scope.
    \item Percentage-based allocation of mixed-load weights or amounts in this version.
    \item Ownership transfer and selectable local-only/cloud-only storage modes.
    \item Interactive maps, drawn property/parcel boundaries, visual area selection
    and map-derived area estimates for the final delivery.
    \item Comparisons between campaigns or harvest years in the final delivery.
    \item Detailed Excel exports in the final delivery.
    \item Dedicated in-application document printing, including batch printing.
    \item An owner-facing business-closure feature, removed from scope under
    \decisionref{Q-SCOPE-07}.
\end{itemize}

Maps, cross-campaign/cross-year comparisons and detailed Excel exports remain
possible future enhancements under \decisionref{Q-SCOPE-02},
\decisionref{Q-SCOPE-03} and \decisionref{Q-SCOPE-04}, respectively. No implementation
date is committed; later inclusion requires an explicit scope-change review.
The final delivery uses property/parcel names, textual location and optional
declared hectares, reviews one campaign at a time and exports authorised campaign
summaries as PDF under \identifier{FR-062}.

Authorised document and generated-PDF downloads are sufficient for the final
delivery under \decisionref{Q-SCOPE-05}. Users may print downloaded files through
their browser or PDF reader. Integrated printing, batch printing and printer
management remain possible future enhancements, with no implementation commitment.

Worker activity is consulted directly in each worker's profile under
\identifier{FR-023}. The final delivery includes days and hours worked, without a
dedicated attendance-document workflow; see \decisionref{Q-SCOPE-06}.

\subsection{Single Final Delivery}
\label{sec:mvp-proposal}
\label{sec:final-delivery}

The project has one final version and one final delivery. No separate MVP release
or staged functional deliveries are planned. Internal development and testing may
proceed progressively towards that same delivery.

The current final-delivery baseline is an end-to-end
harvest workflow in which the owner can:

\begin{enumerate}
    \item Prepare properties, campaign coverage,
    workers and necessary documents.

    \item Plan work and consult prepared schedules.

    \item Confirm attendance and record
    actual periods and tasks in the field,
    including without connectivity.

    \item Record transport and delivery,
    then complete ticket, analysis
    and reported monetary information
    as it becomes available.

    \item Consult relevant insurance,
    use the agreed signing interaction
    and inspect campaign results.
\end{enumerate}

The backlog records implementation order, dependencies, effort, ownership and
acceptance criteria for the final delivery. Any scope reduction or addition needs
an explicit reviewed decision; backlog order alone does not exclude a requirement.
Future enhancements remain outside the final-delivery commitment; the selected four document templates and current-only replacement policy follow the decisions in Section~\ref{sec:validation}. Signing, local document protection and recovery still
require operational validation.

The navigable prototype prepared for the intermediate presentation, described in
Section~\ref{sec:document-status}, is a prior demonstration artefact. It does not
constitute another functional product delivery or acceptance of the final system.

\Needspace{8\baselineskip}
\subsection{Feature Overview}

{\readabletable
    \begin{longtable}{L{0.28}L{0.54}L{0.18}}
\caption{Functional areas and corresponding use cases.}\label{tab:feature-overview}\\
        \toprule
\rowcolor{TableHeader}
        \textbf{Area}
        & \textbf{Main capability}
        & \textbf{Use cases}
        \\
        \midrule

\endfirsthead
\toprule\rowcolor{TableHeader}
\textbf{Area} & \textbf{Main capability} & \textbf{Use cases}\\
\midrule\endhead
\multicolumn{3}{r}{\footnotesize\itshape Continued on next page}\\\endfoot
\bottomrule\endlastfoot
        Account access
        & Owner registration, email verification,
        login, logout and recovery.
        & \identifier{UC-01}--UC-03
        \\
\midrule

        Manager access
        & Invitations, operational defaults, adjustable permissions and revocation.
        & \identifier{UC-04}
        \\
\midrule

        Land and campaign organisation
        & Properties, optional parcels, campaign coverage, planned/actual dates and lifecycle.
        & \identifier{UC-05}--UC-06
        \\
\midrule

        Workforce
        & Identified worker profiles, machinery capability and authorised same-day offline creation.
        & \identifier{UC-07}
        \\
\midrule

        Planning
        & Daily and multi-day plans,
        reuse and schedule consultation.
        & \identifier{UC-08}--UC-09
        \\
\midrule

        Actual work
        & Attendance, task periods,
        breaks and corrections.
        & \identifier{UC-10}--UC-11
        \\
\midrule

        Documents and signing
        & Document access, uploads, retention, offline preparation and worker-specific signing.
        & \identifier{UC-12}--UC-14
        \\
\midrule

        Harvest logistics
        & Transport-associated deliveries,
        weighing tickets and analyses.
        & \identifier{UC-15}--UC-17
        \\
\midrule

        Reported values
        & Current cooperative-reported
        ticket amounts.
        & \identifier{UC-18}
        \\
\midrule

        Campaign review
        & Indicators, supporting detail
        and PDF summaries.
        & \identifier{UC-19}
        \\
\midrule

        Document generation
        & Reviewed template-based generation, followed by version-specific signing.
        & \identifier{UC-20}
        \\

    \end{longtable}
}

Administrative configuration normally requires
connectivity.
Essential field entry and consultation
use prepared authorised information
and explicit synchronisation states.

Offline capability is defined per operation.
It does not imply that the entire application
or every historical record is available
without a connection.

\subsection{Assumptions and Dependencies}

\subsubsection{Business Assumptions}

The baseline provides one business workspace per owner. Several independently
owned workspaces through the same owner account are outside scope. Managers use
owner-assigned property scope, including parcels and current/future campaigns;
several managers may share a property under \decisionref{Q-AA-02}. Business closure
is excluded under \decisionref{Q-SCOPE-07}. The separate procedure for account/data
erasure requests remains under \decisionref{Q-AA-07}; it introduces no closure feature.

The current requirements review assumes that:

\begin{itemize}
    \item Reference information and documents are prepared when possible. New workers
    may be created on the day, including offline, with mandatory identification.
    Their creation does not establish documentary readiness.
    \item A worker may perform different tasks
    and work in several areas during the same day.
    \item Campaigns may cover several properties
    or selected parcels.
    \item Source information may arrive
    after the physical delivery
    or campaign closure.
    \item One transport record is associated
    with each delivery in the final version.
    \item Each delivery may have at most one ticket and one analysis, each added later. Yield aggregation uses the ticket's client-confirmed KILOS ASIGNADOS.
    \item Monetary values are supplied
    by the cooperative and associated
    with weighing tickets.
\end{itemize}

These assumptions must be revised
when validated examples contradict them.

\subsubsection{Business Dependencies}

The team depends on client or intermediary input
to confirm:

\begin{itemize}
    \item Additional source examples, where a new label or exception is not covered by the confirmed ticket interpretation.
    \item Representative assigned weights and analysis values for test fixtures.
    \item Operational review of the four selected Spanish forms and their rendered signing output.
    \item Concrete responsibility assignments and validation of the default permission profile.
    \item Final-delivery scope
    and operational acceptance expectations.
\end{itemize}

Real examples should be provided
with unnecessary personal information removed.

\subsubsection{Technical Dependencies}
\label{sec:technical-dependencies}

Registration, invitations and recovery depend on identity management and email
delivery. Offline access depends on prior authentication, protected local storage
and the authorisation policy in \hyperref[br:002]{\identifier{BR-002}}. Session, password, authentication and logging controls are specified in
Sections~\ref{sec:nfr-security} and~\ref{sec:nfr-offline}. The approved targets
require technical proof; residual mechanisms and deployment choices are listed
in Section~\ref{sec:validation-pending}. Supported financial operations and PDF output use the same grant and
preparation conditions as other field operations.

The solution depends on:

\begin{itemize}
    \item Devices and browsers capable
    of supporting the selected application approach.
    \item Successful preparation of authorised
    reference data before offline work.
    \item Sufficient local storage
    for pending operations and prepared documents.
    \item A central service for validation,
    persistence and synchronisation.
    \item Defined recovery behaviour
    for rejected operations, concurrent changes
    and interrupted transfers.
    \item Consistent deletion propagation
    for document versions removed
    under the retention rules.
\end{itemize}

The selected progressive web application (PWA), IndexedDB, modular backend and PostgreSQL stack is described
in Section~\ref{sec:solution-architecture}. Its selection does not establish
feasibility of the unresolved local-security and synchronisation mechanisms;
these require the technical validation listed in Section~\ref{sec:validation-pending}.

\subsubsection{Project Dependencies}

The academic schedule, team availability,
hosting resources and required course deliverables
constrain the implementation plan.

Use-case approval, task ownership,
effort estimates and acceptance evidence
must be tracked explicitly.
The existence of a documented use case
does not establish that its implementation
has been completed or validated.
The working risk register and response plans in Section~\ref{sec:risk-management}
link these dependencies to task ownership, effort tracking and milestone review.
Final delivery dates, resource commitments and assessment validation remain under
\decisionref{Q-PM-01}.

\subsection{Requirements Review and Change Control}\label{sec:requirements-change}
Treat the current specification as a reviewable snapshot, not permanently frozen
requirements. Record each proposed change with its source, request/review dates,
affected IDs and disposition. Use the use-case links, shared business rules and
architecture/data dependencies to assess scope, effort, schedule, quality and risk
before acceptance.

The owner/client validates business meaning through the intermediary; the team
reviews feasibility and delivery impact. Neither a demonstration nor an unanswered
question counts as approval. Record the decision and reviewer, update affected
artefacts together, and adjust estimates/backlog before committing the change.
Identify the baseline revision used for an estimate or demonstration. The exact
approval assignment and scope-change responsibilities remain under \decisionref{Q-SCOPE-01}
and \decisionref{Q-PM-02}.

\subsection{Operational Document Examples}\label{sec:source-examples}
The client intermediary supplied operational photographs on 6 November 2026 and accompanying
explanations. Their review on 21 November extends the working specification with
concrete source fields and form structures. The examples selected for this specification comprise one cooperative
ticket and four individual worker forms. The following
tables record their interpretation and the client explanations reconciled during this review. The four Spanish forms are the selected generation templates; technical rendering and operational review remain planned under Section~\ref{sec:validation}. Example
names and personal identifiers are omitted. These examples establish no additional
mandatory pre-work document check.

\subsubsection{Cooperative Ticket}
The example identifies cooperative season 2023--2024 and ticket/sample 14011.
Its processing start is 07:08:48, finish 07:14:07 and printed duration 00:05:20.
The timestamps differ by 00:05:19; preserve both printed duration and recorded
times and flag the discrepancy for review. Cooperative processing times are distinct
from trip departure/arrival and worker hours.
{\compacttable\begin{longtable}{L{0.28}L{0.72}}
\caption{Observed ticket fields and their working interpretation.}\label{tab:source-ticket}\\
\toprule\rowcolor{TableHeader}\textbf{Source label} & \textbf{Representation and interpretation}\\\midrule\endfirsthead
\toprule\rowcolor{TableHeader}\textbf{Source label} & \textbf{Representation and interpretation}\\\midrule\endhead
\multicolumn{2}{r}{\footnotesize\itshape Continued on next page}\\\endfoot\bottomrule\endlastfoot
TICKET, FECHA, season & Source ticket reference, local source date and cooperative season label. The source season need not equal the application's campaign name.\\\midrule
H. INI, H. FIN, TIEMPO EMPLEADO & Processing start, finish and independently reported duration, preserving seconds. Validate known chronology and flag discrepancies without overwriting source facts.\\\midrule
LINEA, NUM. MUESTRA & Tractor queue/line at the cooperative and sample reference, according to the supplied explanation. Matching sample and ticket numbers in this example do not imply universal equality.\\\midrule
CALIDAD: VUELO; VARIEDAD & The intermediary identifies VUELO as tree-picked olives and MULTIVARIETAL as mixed varieties. Preserve original labels; other quality codes require confirmation.\\\midrule
TERMINO; POL; PAR; REC; PARAJE & Municipality and printed land identifiers/locality. Preserve literal identifiers as text. The intermediary explains POL as polygon/field number, PAR and REC as field identifiers, and PARAJE as the field name. Preserve these reported labels without requiring a map or automatic cadastral conversion.\\\midrule
SOCIO, member details; MATRICULA; LOCALIZADOR & Reported member identity/reference, vehicle plate when printed, and source locator. These are source snapshots, independent of internal UUIDs and registered worker driver.\\\midrule
KILOS TOTALES; KILOS ASIGNADOS & Total weight and accepted assigned weight, both 929~kg here. The client confirmed that KILOS ASIGNADOS is the cooperative-accepted final weight used for indicators and yield weighting. Store it once as the canonical assigned kg; preserve total kg separately.\\\midrule
MEDIANERIA & Percentage of the total weight accepted by the cooperative, confirmed by the client; 100\% here. It is not laboratory yield. Use reported assigned kg directly, without applying this percentage again or inventing a missing assigned weight.\\\midrule
Penalty, when reported & The intermediary describes a penalty above a time-per-kg allowance and states that this ticket has none. Record any reported penalty, its unit and application status without calculating a time-per-kg charge. No threshold or tariff is required because automatic penalty calculation is outside scope.\\
\end{longtable}}

\subsubsection{Individual Worker Forms}\label{sec:document-examples}
The four individual forms use company, worker name, DNI/NIE, job role, date and
worker signature. Job role is a document snapshot, separate from a work period's
actual task. The supported identification vocabulary follows \decisionref{Q-ID-01}. Generated documents retain the supplied Spanish wording, while the application interface is English.
{\compacttable\begin{longtable}{L{0.3}L{0.7}}
\caption{Supplied worker form inventory.}\label{tab:source-forms}\\
\toprule\rowcolor{TableHeader}\textbf{Form} & \textbf{Specific content and treatment}\\\midrule\endfirsthead
\toprule\rowcolor{TableHeader}\textbf{Form} & \textbf{Specific content and treatment}\\\midrule\endhead
\multicolumn{2}{r}{\footnotesize\itshape Continued on next page}\\\endfoot\bottomrule\endlastfoot
Personal protective equipment receipt & Reviewed equipment selections, optional other-item descriptions and acknowledgement of delivery/use information. Store selections in the generated source snapshot.\\\midrule
Safety information receipt & Reviewed list of supplied risk/prevention information, equipment/product instructions and emergency information. Signature records the acknowledgement represented by the source.\\\midrule
Machinery-use authorisation & Reviewed list of authorised machinery. Store the list in the document; profile machinery capability and actual work tasks remain separate facts.\\\midrule
Medical-examination consent & Explicit acceptance or refusal, date and worker signature in the corresponding area. A signed refusal is a completed signing result with refused consent. Clinical results are outside this form workflow.\\
\end{longtable}}
The four individual forms are classified as \textit{other}, with explicit form
subtypes. Every document belongs to exactly one worker. Generating for several workers
produces separate individual documents/sources from the selected template.
Insurance-only manager access does not include these forms. Template wording and
legal references require client review before operational use. PPE stock management, automated machinery
authorisation enforcement and clinical records are not introduced by these examples.


\paragraph{Private access and retained functional scope.}
Owner self-registration and onboarding additional independent client businesses
are outside the final delivery. Technical administration covers initial setup and
activation support only. Owner-managed invitations, property-scoped manager grants,
optional planning/schedule consultation and the remaining agricultural workflows
retain their specified behaviour. The administrator is not an operational substitute
for the owner or manager. Handover includes account activation, representative
phone/tablet preparation, a short user guide and named technical support ownership.
